\documentclass[11pt,a4paper]{article}
% One-page summary of the master's thesis (for the examiners, colloquium email).
\usepackage[margin=22mm,top=20mm,bottom=20mm]{geometry}
\usepackage[T1]{fontenc}
\usepackage{mathptmx}
\usepackage{enumitem}
\usepackage{xcolor}
\definecolor{c3}{RGB}{30,80,160}
\setlength{\parindent}{0pt}
\setlength{\parskip}{3pt}
\setlist{nosep,leftmargin=1.2em}
\pagestyle{empty}
\newcommand{\sect}[1]{\par\vspace{5pt}{\color{c3}\bfseries #1}\par\vspace{1pt}}

\begin{document}
{\large\bfseries Hamilton-Principle Models of Oral Biofilm Dysbiosis:
GPU-Accelerated Inference, Metabolic Priors, and Spatial Stratification}\par
\vspace{3pt}
Master's thesis, Keisuke Nishioka, Institut für Kontinuumsmechanik, Leibniz
Universität Hannover (double degree with Keio University).\\
Examiners: Prof.\ Dr.-Ing.\ Philipp Junker, Dr.-Ing.\ Matthias Wangenheim.
Supervisors: Dr.-Ing.\ Meisam Soleimani, Dr.-Ing.\ Hendrik Geisler.
\par\vspace{2pt}\hrule

\sect{Motivation}
Peri-implantitis is driven by multi-species biofilms whose composition shifts
from commensal to dysbiotic. A point model of a multi-species biofilm (Klempt
et al.\ 2026) describes how the composition at one material point evolves, but
has no space. A continuum growth model (Klempt et al.\ 2024) describes where a
biofilm grows and which stresses it carries, but does not distinguish species.
Both follow from the extended Hamilton principle. The task of this thesis is to
bring the point model into the existing finite element framework of the IKM, so
that composition, growth and stress can be computed together.

\sect{Contributions}
\begin{enumerate}
\item \textbf{A calibrated point model.} The 15 interaction coefficients of a
  five-species point model are estimated with transitional Markov chain Monte
  Carlo from the in-vitro time courses of Heine et al.\ (2025) under four
  conditions. Compiled with JAX, the forward model runs about 150 times faster
  on the GPU than on the CPU. The fits follow all four conditions (RMSE 0.03 to
  0.12), and the pH, which was not used in the calibration, is predicted with
  $R^2=0.78$. A metabolic sign prior carries the same type of model to sparse
  in-vivo data of ten patients.
\item \textbf{A growth law in the finite element framework.} The growth
  kinematics $\mathbf F=\mathbf F_e\mathbf F_g$, $\mathbf F_g=\alpha\mathbf I$,
  is implemented as an ANSYS user material, verified against closed-form
  solutions, and coupled to the biofilm field of the IKM element through the
  growth equation $\dot\alpha=k_\alpha\phi$ at every Gauss point. In ANSYS the
  coupled run follows the exact solution of a gradient-free region to
  $3\times10^{-11}$ in $\alpha$.
\item \textbf{The composition at every Gauss point.} For two species, the point
  model runs at every Gauss point through a bridge from Fortran to Python: the
  field sets the amount of biofilm, the point model its composition. Every call
  is reproduced in Python bit for bit, and the published outcomes, coexistence
  and the survival of one species, are kept.
\end{enumerate}

\sect{Towards a two-way coupling}
Computed in Python on the reproduced field of Klempt et al.\ (2024): with the
local nutrient in the point model the composition varies across the biofilm, and
with a composition-dependent front rate, as in Soleimani et al.\ (2023) and Feng
et al.\ (2021), the amount of biofilm changes by 4 to 13\,\%, depending on the rate contrast.
The first of these steps also runs in ANSYS: with the element's nutrient passed
to the point model, the composition in the seed varies in space. Measured
expansion data of Fei et al.\ (2020) indicate that the growth rate must depend on
the local nutrient.

\sect{Limitations and outlook}
All finite element results are verified, not validated: no measured stress of a
growing oral biofilm is available. The size of the stresses depends on the
diffusion coefficient of the field, whose conversion to the model is an
assumption; with the converted value the stresses converge in the mesh. In
ANSYS only the local nutrient acts back on the point model, and the front term
of the element is inactive, so the biofilm spreads only by diffusion. More
species, growth towards the nutrient and a validation of the mechanics continue
at Keio University.
\end{document}
